\begin{banner}
REMEDIËRING WISKUNDE \quad·\quad 6 OPVOEDING EN BEGELEIDING
\end{banner}

\begin{titel}
Statistiek, deel 1
\end{titel}

\begin{ondertitel}
Basisbegrippen: van onderzoeksvraag tot kritisch besluit
\end{ondertitel}

Elke dag lees of hoor je cijfers: ‘9 op 10 ouders is tevreden’, ‘jongeren slapen te weinig’, ‘dit product werkt bewezen’. Als begeleider, opvoeder of zorgverlener moet je kunnen beoordelen of zo’n uitspraak klopt. In deze bundel leer je hoe een statistisch onderzoek verloopt, hoe je een \textbf{betrouwbare steekproef} samenstelt, hoe cijfers en grafieken je kunnen \textbf{misleiden} en waarom \textbf{samenhang} nog geen \textbf{oorzaak} is.

Je werkt je door zes levels, zoals in een spel. Na elk level test je jezelf met een \textbf{CHECK}. Haal je de score? Dan ga je verder. Lukt het nog niet? Herneem dan de theorie en de oefeningen. Op het einde wacht de \textbf{eindtoets}.

\textbf{Naam:} \_\_\_\_\_\_\_\_\_\_\_\_\_\_\_\_\_\_\_\_\_\_\_\_\_\_\_\_\_\_ \quad \textbf{Klas:} \_\_\_\_\_\_\_\_\_\_\_\_\_\_ \quad \textbf{Datum:} \_\_\_\_\_\_\_\_\_\_\_\_\_\_

\subsection*{Jouw route}

\breedtes{12/50/18/12}
\begin{tabular}{llll}
\hline
Level & Onderwerp & Leerplandoel & Behaald? \\
\hline
1 & Statistisch onderzoek: stappen en variabelen & basis (LPD 1, 3) & ☐ \\
2 & Steekproef en populatie & LPD 11 & ☐ \\
3 & Steekproefontwerp: representatief of niet? & LPD 11 & ☐ \\
4 & Misleidingen en valkuilen & LPD 1 – 3 & ☐ \\
5 & Samenhang versus causaliteit & LPD 12 & ☐ \\
6 & Eindmissie: onderzoeken, problemen oplossen, uitleggen & LPD 1 – 4 & ☐ \\
\textbf{Eind} & \textbf{Eindtoets} & alle doelen van deze bundel & ☐ \\
\hline
\end{tabular}

\pagina

\section{Zo werk je met deze bundel}

\begin{missie}
Missie: weten wat je moet doen, en hoe je het doet.
\end{missie}

\begin{theorie}
\textbf{THEORIE} \quad De leerstof, steeds met een voorbeeld. Lees ze rustig en werk mee met potlood en rekenmachine.
\end{theorie}

\begin{voorbeeld}
\textbf{VOORBEELD} \quad Een uitgewerkt voorbeeld uit de wereld van opvoeding, zorg en welzijn.
\end{voorbeeld}

\begin{letop}
\textbf{LET OP} \quad Hier gaan leerlingen vaak de mist in. Lees deze kaders extra goed.
\end{letop}

\begin{ict}
\textbf{ICT-TIP} \quad Zo laat je Excel, GeoGebra of je rekentoestel het rekenwerk doen.
\end{ict}

\begin{doel}
\textbf{DOEL} \quad Dit moet je kunnen na het level. Kruis het vakje aan wanneer het lukt.
\end{doel}

\subsection*{De oefeningen}

Elke oefening heeft een moeilijkheidsgraad: \textbf{★☆☆} (opwarmer), \textbf{★★☆} (gewone oefening) en \textbf{★★★} (uitdaging). Je vindt ook steeds deze soorten opdrachten:

\breedtes{22/78}
\begin{tabular}{ll}
\hline
Soort & Wat moet je doen? \\
\hline
\textbf{Aan de slag} & Je past de theorie toe op een situatie of op getallen. \\
\textbf{Context} & Een situatie uit de opvoedings-, zorg- of welzijnswereld. Je zoekt zelf welk begrip je nodig hebt. \\
\textbf{Invullen} & Een tekst met ontbrekende woorden. Kies uit de woordenlijst. \\
\textbf{Koppel} & Je koppelt een begrip aan een situatie, een grafiek of een omschrijving. \\
\textbf{Foutjager} & Iemand maakte een fout. Zoek ze, leg uit en verbeter. \\
\textbf{CHECK} & Mini-toets op het einde van elk level. \\
\hline
\end{tabular}

\begin{letop}
\textbf{Schrijf je werkwijze op en eindig met een antwoordzin.} Zo toon je dat je het resultaat ook begrijpt.
\end{letop}

\subsection*{Doelenkaart: wat moet je kunnen?}

Vul in \textbf{voor} je begint en opnieuw \textbf{na} de eindtoets: \textbf{+} (lukt zeker), \textbf{±} (lukt half), \textbf{–} (lukt nog niet).

\breedtes{8/68/12/12}
\begin{tabular}{llll}
\hline
LPD & Ik kan… & Voor & Na \\
\hline
 & de stappen van een statistisch onderzoek ordenen, variabelen en soorten data herkennen. & & \\
11 & het verschil tussen populatie en steekproef uitleggen, redenen voor een steekproef noemen en steekproefvariabiliteit verklaren. & & \\
11 & verklaren waarom een steekproef representatief moet zijn, de criteria toepassen en vertekende steekproeven (opportunistisch, vrijwillige respons, non-respons) herkennen. & & \\
12 & in concrete situaties het verschil uitleggen tussen samenhang en causaliteit en observationeel en experimenteel onderzoek onderscheiden. & & \\
1 – 4 & misleidende grafieken, cijfers en uitspraken herkennen (1), een probleem aanpakken met een antwoordzin (2), ICT gebruiken (3) en mijn redenering uitleggen (4). & & \\
\hline
\end{tabular}

\begin{letop}
\textbf{Opmerking} \quad Deze bundel behandelt de leerplandoelen over steekproeven en samenhang (LPD 11 en 12) en de doorlopende doelen LPD 1 – 4. De doelen over de normale verdeling (LPD 13 en 14) komen aan bod in de bundel ‘Statistiek 6OB’ (deel 2: kengetallen en normale verdeling).
\end{letop}
